% ============================================================================
% CONCLUSION DU CHAPITRE 1 
% source : partie1/conclusion_partie1.tex
% ============================================================================

\section{CONCLUSION DU CHAPITRE~\ref{chap:contexte}}
\label{sec:conclusion_contexte}

\begin{sloppypar}
Le chapitre~\ref{chap:contexte} a pos\'e progressivement les fondements sur lesquels
s'appuie l'ensemble de la th\`ese. Il en ressort trois acquis.

Le premier acquis, de nature illustrative, porte sur trois syst\`emes cryptographiques d\'eploy\'es,
\`a savoir Helios pour le vote \'electronique, les zk-SNARKs pour la conformit\'e
financi\`ere et ECDSA pour la signature contractuelle. Selon des
estimations param\'etriques sur l'\'echelle op\'erationnelle CLO, leur produit
$\Conf\cdot\Relver\cdot\Opp$ est compris entre $0{,}273$ et $0{,}572$. Ce motif
laisse entrevoir, sans la d\'emontrer, une tension structurelle entre deux paradigmes
de v\'erification.
Le paradigme cryptographique mesure la r\'esistance aux attaques, alors que le
paradigme institutionnel requiert l'intelligibilit\'e par des non-experts,
des traces exploitables pour un audit r\'etrospectif ainsi qu'une responsabilit\'e
attribuable en cas de d\'efaillance. Aucune m\'etrique formelle n'avait encore
\'etabli les limites de ce compromis.

Le deuxi\`eme acquis est conceptuel. L'opposabilit\'e
$\Opp = I(V;\Sigma|C)/H(V|C)$ entre comme troisi\`eme dimension
cryptographique aux c\^ot\'es de la confidentialit\'e et de la
fiabilit\'e. Deux arch\'etypes s'en d\'egagent, les syst\`emes \`a
haute confidentialit\'e (Helios, zk-SNARKs, produits $0{,}437$ et
$0{,}572$) s'opposant aux syst\`emes \`a haute opposabilit\'e (ECDSA,
produit $0{,}273$). L'\'equilibre hypoth\'etique
$\Conf \approx \Rel \approx \Opp \approx 0{,}8$ n'est approch\'e par aucune
des trois estimations. Cette d\'efinition transforme
l'interpr\'etabilit\'e institutionnelle d'une contrainte externe
subie en une propri\'et\'e formelle que l'on peut mesurer et optimiser.

Le troisi\`eme est programmatique. La cryptographie th\'eorique
formalise $\Conf$ et $\Rel$ via des jeux de s\'ecurit\'e mais ignore
$\Opp$~; la s\'ecurit\'e utilisable analyse l'utilisabilit\'e
qualitativement sans formalisation math\'ematique~; l'informatique
juridique raisonne sur l'admissibilit\'e de la preuve num\'erique
sans m\'etriques quantitatives. Les six objectifs de recherche
Obj~\ref{obj:obj1_unification} \`a Obj~\ref{obj:obj6_experimentation} d\'eclinent
le paradoxe de l'authenticit\'e invisible en probl\`emes formels
pr\'ecis, chacun associ\'e \`a une contribution et \`a une \'etape
de la progression canonique
\textbf{F}$\to$\textbf{B}$\to$\textbf{E}$\to$\textbf{S}$\to$\textbf{C}$\to$\textbf{X},
comblant ainsi le cloisonnement entre ces trois corpus.

Le chapitre~\ref{chap:cadre_formel} prend le relais. Il dispose d\'esormais
d'un vocabulaire fix\'e, d'une probl\'ematique formul\'ee avec pr\'ecision
et d'exemples calibr\'es. La mesure de $\Conf$, $\Rel$ et $\Opp$ remplace le d\'ebat qualitatif entre s\'ecurit\'e et transparence et transforme des choix de conception en
probl\`emes d'optimisation sous contraintes, susceptibles de r\'eponses
pr\'ecises. Il peut d\'esormais engager les six \'etapes du programme,
de la formalisation \`a l'exp\'erimentation.
\end{sloppypar}

